\section{Core language: dependence graph semantics}

\figref{graph:core:eval} defines a small-step reduction relation.

Notes:
\begin{enumerate}
\item No support for recursion yet.
\item If ``reduct'' (source of hyperedge) means new graph created on right by the rule, then both projection and \kw{let} rules result in an empty reduct (hyperedge with an empty source). No path can connect to such a hyperedge, so backward slicing will never reach such a redex. For forward slicing to be adjoint to that behaviour, redexes with empty reducts need to be ``costless'', i.e.~always included by forward slicing.
\end{enumerate}

\subsection{Foreign functions}

Foreign functions must be defined precisely with inputs that are appropriately truncated. For example, a binary operator $*$ with annihilator $0$ must define rules such as:

\begin{salign}
\annot{\exInt{0}}{\alpha} * \beta &\geval \annot{\exInt{0}}{\alpha'} \\
\annot{\exInt{n}}{\alpha} * \annot{\exInt{n}}{\beta} &\geval \annot{(\exInt{n \mathbin{\hat{*}} m})}{\alpha'} & \textit{if }n \neq 0 \\
\end{salign}

\begin{figure}
   {\small \flushleft \shadebox{$r \geval r'$}}%
   \begin{salign}
      \annRecProj{\annRec{(\bind{x}{v}) \cons \seq{\bind{y}{u}}}{\beta}}{x}{\alpha}
      & \geval
      v
      &
      \\
      \annApp{\annFun{\sigma}{\beta}}{v}{\alpha}
      & \geval
      \appSubst{\gamma}{e}
      &
      \text{if }v, \sigma \match \gamma, e
      \\
      \annLet{x}{v}{e}{\alpha}
      & \geval
      \appSubst{(\bind{x}{v})}{e}
      \\
      \annForeignApp{f}{\seq{v}}{\alpha}
      & \geval
      \interpret{f}(\seq{v})
   \end{salign}
   \caption{Core language: small-step reduction relation}
   \label{fig:graph:core:eval}
\end{figure}

